\section{本章小结}
\label{sec:opt-prob-summary}

章首把优化想象成在一块允许站的地上找最低点。本章把这幅画面拆成了几个分别回答的问题。问题的表述要同时给出决策变量、目标函数与可行域；最优值是下确界，最优解是取得它的可行点，二者可以分离：问题可能不可行、无下界，或最优值有限却未达到。局部最优只与附近的可行点比较，全局最优与所有可行点比较；扩展实值写法用指示函数把约束并入目标，不改变最优值与最优解，而有限惩罚会改变问题。

两个贯穿例子展示了这些问题怎样得到回答。资源分配的可行域是有界多面体，紧集上的连续目标保证最优解存在；把两条资源约束相加，就同时得到利润上界$10$和达到它的方案$(2,2)$。最小二乘的可行域无界，正则项$\lambda>0$使目标强制，从而最优解存在；同一个正则项使Hessian$\bm A^\top\bm A+\lambda\bm I$正定，目标严格凸乃至$\lambda$强凸，于是最优解唯一，目标差还能换成到最优解的距离。$\lambda=0$且$\bm A$秩亏时，这些结论依次失效：目标不强制，最优解仍存在却成一条直线，参数不唯一而给定样本上的预测相同。

表~\ref{tab:opt-prob-summary}汇总了各个结论所需的条件。表中的条件都是充分条件，反例列说明去掉条件或反向推断时会出什么问题。

\begin{table*}[!htbp]
\centering
\caption{本章主要问题、充分条件与边界。表中“凸”均指可行域凸且目标凸；强凸常数$\mu$对应欧氏范数。}
\label{tab:opt-prob-summary}
\small
\begin{tabularx}{\linewidth}{@{}>{\RaggedRight\arraybackslash}p{24mm}>{\RaggedRight\arraybackslash}X>{\RaggedRight\arraybackslash}X>{\RaggedRight\arraybackslash}X@{}}
\toprule
问题 & 充分条件 & 结论 & 反例或边界 \\
\midrule
最优解是否存在 & 可行域非空紧，目标适当且下半连续；或可行域非空闭，目标适当、下半连续且强制 & 最优值有限且被达到 & $(0,1]$上的$x$；$[1,\infty)$上的$1/x$；在$0$处跳高的函数；无限维中有界闭集不紧 \\
局部是否即全局 & 凸 & 每个局部最优解都是全局最优解，最优解集合凸 & $x^4/4-x^3/3-x^2$在$x=-1$处的局部最优 \\
最优解是否唯一 & 凸可行域上严格凸，且最优解存在 & 最优解唯一 & $e^x$严格凸而无最小点；$|x|$与资源分配问题有唯一解而目标不严格凸 \\
目标差能否换成距离 & 凸可行域上$\mu$强凸，且最优解存在 & $f(\bm x)-f^\star\geq\frac{\mu}{2}\lVert\bm x-\bm x^\star\rVert^2$ & 非凸可行域$\{|x|\geq1\}$上的$x^2$ \\
内点能否用梯度判别 & 可微；凸目标在开凸集上 & 局部最优必有$\nabla f=\bm 0$；凸时$\nabla f=\bm 0$即全局最优 & 鞍点$x^2-y^2$；边界最优点导数可不为零 \\
驻点是否局部最优 & 二阶连续可微，Hessian正定 & 严格局部最优解 & $x^2+y^4$（正定非必要）；$x^2+y^3$（半正定不充分） \\
小残差是否表示小误差 & 可微且$\mu$强凸 & $f(\hat{\bm x})-f^\star\leq\lVert\nabla f(\hat{\bm x})\rVert_2^2/(2\mu)$ & 平坦目标$0.01x^2/2$；$\lambda=0.01$的最小二乘 \\
\bottomrule
\end{tabularx}
\end{table*}

最后一节说明了怎样评价实际交出的近似解。目标差、到解集的距离与可计算的残差是三种不同的误差，只有在强凸这类附加条件下才能互相换算，换算系数随$1/\mu$放大。比较方法时，还要把达到指定精度的轮数与每轮代价合成为总成本，并在相同的问题、误差对象、精度和规模下比较。

本章没有生成任何迭代序列，也没有假设任何算法已经成功。我们得到的是后续各章共同使用的语言：有限维的凸几何、解的存在与唯一、内点的最优条件，以及评价近似解的口径。下一章在最具体的线性模型上把这些语言用于一个完整的求解过程：可行域是多面体，最优解在顶点中寻找，而“把约束加权相加得到下界”的想法将发展为可以核对的最优性证书。
